\section{Evaluación de modelos}
\label{sec:marco_evaluacion}

De los tres tipos de salida, las decisiones de abstención ya tienen su medida: la cobertura y el riesgo selectivo
de la Ecuación~\ref{eq:marco_riesgo_selectivo}, calculados con la métrica de la salida que se acepta. Esta sección
presenta las métricas de los otros dos, las categorías y los textos, y las propiedades que se evalúan en ambos.
Antes, sin embargo, conviene establecer qué puede afirmar una métrica calculada sobre el conjunto de prueba
(\emph{test}), porque de eso depende cómo se leen todas las demás.

\subsection{Sesgo y generalización}
\label{sec:marco_generalizacion}

Una métrica calculada sobre el test estima el desempeño del modelo en la distribución de la que proviene ese test,
que suele ser la misma fuente que la del entrenamiento. En uso real, las imágenes llegan de otros equipos, centros y
pacientes, y el cambio de distribución de la Sección~\ref{sec:marco_regularizacion} reduce el desempeño de forma
sustancial en problemas de todo tipo, aun con los métodos diseñados para mitigarlo. Ese cambio toma dos formas: el
\textbf{cambio de dominio} (\emph{domain shift}), cuando las imágenes provienen de una fuente que el entrenamiento no vio, y el
\textbf{cambio de subpoblación} (\emph{subpopulation shift}), cuando la fuente es la misma pero los grupos que la componen aparecen en otras
proporciones~\cite{koh2021wilds}.

En dermatología, el cambio de dominio más visible es el de modalidad. La fotografía clínica registra la lesión tal
como se ve a simple vista, y la dermatoscopía la examina con aumento e iluminación controlada, de modo que muestra
estructuras que la fotografía no resuelve, como la red de pigmento y los patrones
vasculares~\cite{dinnes2018dermoscopy,kittler2016standardization}. Un modelo entrenado con una modalidad y usado
con la otra recibe imágenes cuyos rasgos visuales se distribuyen de otra manera que en el entrenamiento.

Aun dentro de una misma modalidad, la iluminación, la distancia a la lesión, la resolución del sensor y la
compresión de la imagen varían entre equipos y centros. Un modelo que se apoya en rasgos que esas variaciones
alteran, como la textura fina o el color exacto, puede perder desempeño ante imágenes de contenido clínico
idéntico. Esa pérdida se estima de dos maneras. La \textbf{validación externa} (\emph{external validation}) mide
el modelo en un conjunto de otra fuente, que no intervino en el entrenamiento. Cuando no se dispone de otro
conjunto, se degradan las imágenes del test de manera controlada (Sección~\ref{sec:marco_eval_robustez}).

El tono de piel es el ejemplo más estudiado de cambio de subpoblación. Cuando las pieles claras están mucho más
representadas que las oscuras en el entrenamiento, el modelo clasifica mejor los fototipos parecidos a los que
vio~\cite{groh2021neural}. Evaluados sobre imágenes con fototipos diversos, modelos dermatológicos del estado del
arte rindieron peor que en sus conjuntos de origen, sobre todo en pieles oscuras y en enfermedades poco
comunes~\cite{daneshjou2022disparities}. Una diferencia sistemática de desempeño entre grupos de pacientes es un
\textbf{sesgo} del modelo, y una métrica agregada puede ocultarla cuando el grupo afectado es una fracción pequeña
del test. Por eso la métrica se calcula también por grupo, como el fototipo o la modalidad, en una
\textbf{evaluación desagregada} (\emph{disaggregated evaluation}).

Un tercer problema no aparece en ninguna métrica calculada sobre el test. La
Sección~\ref{sec:marco_regularizacion} mencionó que una red puede aprender una regla o una marca de tinta que
aparece junto a ciertas lesiones. Si en un conjunto esas marcas, las etiquetas adhesivas o los bordes oscuros del
dermatoscopio acompañan con más frecuencia a los melanomas, el modelo puede predecir melanoma cuando los ve y
acertar sin mirar la lesión. A estos rasgos, que acompañan a una clase sin formar parte de ella, se los llama
\textbf{rasgos espurios} (\emph{spurious features}), y su efecto se documentó en conjuntos públicos de dermatoscopía~\cite{yan2023towards}.
Los propios modelos contrastivos médicos permiten buscarlos en un conjunto antes de entrenar sobre
él~\cite{kim2024transparent}. Como el rasgo está también en el test, el modelo acierta ahí por la misma razón
equivocada, y para advertirlo hace falta ver en qué parte de la imagen se apoya
(Sección~\ref{sec:marco_eval_gradcam}).

Una métrica sobre el test describe, entonces, el desempeño en una fuente y en una mezcla de grupos, y no dice en
qué se apoya el modelo. Las subsecciones que siguen presentan primero las métricas de cada tipo de salida y después
las propiedades que se evalúan en cualquiera de ellas, entre las que están la robustez y la interpretabilidad, que
estiman lo que la métrica sobre el test no muestra.

\subsection{Evaluación de la clasificación}
\label{sec:marco_eval_clasif}

Evaluar un clasificador es averiguar en qué clases acierta y con cuáles las confunde. Las clases están
desbalanceadas (Sección~\ref{sec:marco_desbalance}) y no todos los errores tienen la misma gravedad, porque no
detectar un melanoma es más grave que derivar un nevo (Sección~\ref{sec:marco_incertidumbre}). Por eso una sola
cifra no alcanza: un modelo puede acertar casi todas las imágenes de la clase mayoritaria y, aun así, clasificar
de forma errónea la mayoría de los melanomas. La herramienta que describe con más detalle los aciertos y los
errores de un clasificador es la \textbf{matriz de confusión}
$M \in \mathbb{N}^{C \times C}$, en la que $M_{ij}$ cuenta las imágenes de la clase real $i$ que el modelo predice
como $j$, y de la que salen todas las demás métricas. En la tarea binaria, con una clase de interés (la positiva)
frente al resto, tiene cuatro celdas: los verdaderos positivos ($\mathrm{TP}$), los falsos negativos
($\mathrm{FN}$), los falsos positivos ($\mathrm{FP}$) y los verdaderos negativos ($\mathrm{TN}$). Cada fila se lee
contra la clase real y cada columna contra la predicha (Figura~\ref{fig:marco_matriz_roc}a). Las métricas
que se derivan de ella son las siguientes.

% Figura: matriz de confusión binaria y curva ROC ilustrativas, con datos sintéticos (no son resultados del
% trabajo). Se usa en sec:marco_eval_clasif, al presentar la matriz de confusión y las métricas que salen de ella.

\begin{figure}[H]
    \centering
    \begin{tikzpicture}[x=1mm, y=1mm]

        % (a) Matriz de confusión binaria.
        \begin{scope}
            \foreach \x/\y/\txt/\col in {0/20/{TP}/figverde, 20/20/{FN}/figbermellon,
                                         0/0/{FP}/figbermellon, 20/0/{TN}/figverde} {
                \fill[\col!15] (\x, \y) rectangle ++(20, 20);
                \draw[figtrazo] (\x, \y) rectangle ++(20, 20);
                \node[font=\small\bfseries] at (\x + 10, \y + 10) {\txt};
            }
            \node[font=\footnotesize] at (10, 43) {positiva};
            \node[font=\footnotesize] at (30, 43) {negativa};
            \node[font=\small] at (20, 49) {clase predicha};
            \node[font=\footnotesize, anchor=east] at (-1.5, 30) {positiva};
            \node[font=\footnotesize, anchor=east] at (-1.5, 10) {negativa};
            \node[font=\small, rotate=90] at (-19, 20) {clase real};

            % Lectura por fila (sensibilidad, especificidad) y por columna (precisión).
            \draw[figflecha] (41, 30) -- (47, 30)
                node[right, font=\footnotesize, align=left] {sensibilidad\\$\mathrm{TP}/(\mathrm{TP}+\mathrm{FN})$};
            \draw[figflecha] (41, 10) -- (47, 10)
                node[right, font=\footnotesize, align=left] {especificidad\\$\mathrm{TN}/(\mathrm{TN}+\mathrm{FP})$};
            \draw[figflecha] (10, -1) -- (10, -7)
                node[below, font=\footnotesize, align=center] {precisión\\$\mathrm{TP}/(\mathrm{TP}+\mathrm{FP})$};
            \node[font=\small] at (20, -24) {(a) Matriz de confusión};
        \end{scope}

        % (b) Curva ROC.
        \begin{scope}[xshift=95mm]
            \fill[figazul!15] (0, 0) -- (0, 22) .. controls (4, 35) and (12, 39) .. (40, 40) -- (40, 0) -- cycle;
            \draw[figtrazo] (0, 0) rectangle (40, 40);
            \draw[figtrazo, dashed] (0, 0) -- (40, 40);
            \draw[figazul!80!black, thick] (0, 0) -- (0, 22) .. controls (4, 35) and (12, 39) .. (40, 40);
            \fill[fignaranja!80!black] (6.2, 32.4) circle (1.1);
            \node[figetiqueta, anchor=west] at (8, 28) {umbral $\tau$};
            \node[font=\footnotesize] at (27, 14) {AUC};
            \node[font=\footnotesize, rotate=38] at (22, 18.5) {azar};
            \foreach \t in {0, 0.5, 1} {
                \node[below, font=\footnotesize] at (\t*40, 0) {\t};
                \node[left, font=\footnotesize] at (0, \t*40) {\t};
            }
            \node[font=\small] at (20, -7) {FPR $= 1 -{}$especificidad};
            \node[font=\small, rotate=90] at (-8, 20) {TPR $={}$sensibilidad};
            \node[font=\small] at (20, -24) {(b) Curva ROC};
        \end{scope}

    \end{tikzpicture}
    \caption{Matriz de confusión binaria con las métricas que se leen en ella, y curva ROC con su área.}
    \label{fig:marco_matriz_roc}
\end{figure}


\begin{itemize}
    \item \textbf{Exactitud} (\emph{accuracy}). Es la fracción de imágenes bien clasificadas, $\sum_c M_{cc} / n$.
    Es engañosa
    cuando las clases están desbalanceadas: si una clase reúne dos tercios de las imágenes, un modelo que responda
    siempre esa clase alcanza una exactitud cercana a 0.67.

    \item \textbf{Sensibilidad macro} (\emph{sensitivity}). La sensibilidad de la clase $c$,
    $R_c = M_{cc} / \sum_{j} M_{cj}$, es la
    fracción de sus imágenes que el modelo acierta, y se lee por fila. La sensibilidad macro promedia esas
    fracciones para que cada clase pese lo mismo, y corrige así el defecto de la exactitud,
    \begin{equation}
        \text{sensibilidad macro} = \frac{1}{C}\sum_{c=1}^{C} R_c.
        \label{eq:marco_acc_bal}
    \end{equation}
    En el caso binario se la llama exactitud balanceada~\cite{brodersen2010balanced,sokolova2009systematic}.

    \item \textbf{Precisión} (\emph{precision}) \textbf{y F1.} La sensibilidad no dice cuántas de las predicciones
    de una clase son correctas. Esa es la precisión de la clase $c$, $P_c$, que se lee por columna. El F1 las
    combina mediante su media armónica, y su promedio sobre las clases se puede hacer de dos maneras,
    \begin{equation}
        F1_c = \frac{2\,P_c R_c}{P_c + R_c}, \qquad
        F1_{\text{macro}} = \frac{1}{C}\sum_{c=1}^{C} F1_c, \qquad
        F1_{\text{ponderado}} = \sum_{c=1}^{C} \frac{n_c}{n}\, F1_c,
        \label{eq:marco_f1}
    \end{equation}
    donde $n_c$ es el número de imágenes de la clase $c$ y $n$ el total. El F1 ponderado queda dominado por la
    clase más frecuente, y el F1 macro da a las $C$ clases el mismo peso.

    \item \textbf{Sensibilidad y especificidad} (\emph{specificity}) \textbf{de una clase.} Cuando interesa una
    sola condición, la salida
    multiclase se colapsa en una binaria: la clase de interés, o un grupo de clases como las malignas, frente a
    las demás, con la probabilidad de la clase o la suma de las del grupo como puntaje. A esa tarea derivada se
    la llama detección, y se evalúa con la sensibilidad $R_c$ y con la especificidad, la fracción de las imágenes
    de las otras clases que el modelo no predice como $c$,
    \begin{equation}
        \mathrm{Esp}_c = \frac{\sum_{i \neq c}\sum_{j \neq c} M_{ij}}{\sum_{i \neq c}\sum_{j} M_{ij}}.
        \label{eq:marco_sens_esp}
    \end{equation}
    Las dos se mueven en sentidos opuestos cuando cambia el umbral de decisión~\cite{fawcett2006introduction}.

    \item \textbf{Curva ROC y \gls{auc}.} El modelo asigna a cada imagen un puntaje $s$ y la declara positiva si
    $s$ supera un umbral $\tau$. Cada umbral produce una tasa de verdaderos positivos, que es la sensibilidad, y
    una de falsos positivos, que es uno menos la especificidad,
    \begin{equation}
        \mathrm{TPR}(\tau) = \frac{\mathrm{TP}}{\mathrm{TP} + \mathrm{FN}}, \qquad
        \mathrm{FPR}(\tau) = \frac{\mathrm{FP}}{\mathrm{FP} + \mathrm{TN}}.
        \label{eq:marco_tpr_fpr}
    \end{equation}
    La curva ROC recorre todos los umbrales y dibuja la TPR contra la FPR (Figura~\ref{fig:marco_matriz_roc}b),
    y el \gls{auc} es el área bajo esa curva. No depende de un umbral, vale 0.5 para un clasificador al azar y 1
    para uno que separa perfectamente las dos clases, y es la probabilidad de que el modelo asigne más puntaje a
    un positivo elegido al azar que a un negativo elegido al azar,
    \begin{equation}
        \mathrm{AUC} = P\!\left(s(x^{+}) > s(x^{-})\right).
        \label{eq:marco_auc}
    \end{equation}
    Se calcula con la regla del trapecio. En el caso multiclase se construye una curva por clase, que trata esa
    clase como positiva y al resto como negativa~\cite{fawcett2006introduction}, y el promedio simple de las $C$
    áreas es el AUC macro, $\mathrm{AUC}_{\text{macro}} = \frac{1}{C}\sum_{c} \mathrm{AUC}_c$.
\end{itemize}

Estas métricas evalúan una predicción que es una categoría: indican en qué clases acierta el modelo, si
descuida las minoritarias y qué tan bien separa cada clase del resto. Por otro lado, un reporte es un texto, y
compararlo con otro exige métricas distintas.

\subsection{Evaluación de la generación de reportes}
\label{sec:marco_eval_texto}

Un reporte se juzga por los dos requisitos de la Sección~\ref{sec:marco_rep_ejes}: que describa lo que la imagen
muestra y que nombre el diagnóstico correcto. Medirlos es más difícil que medir una categoría, porque una misma
imagen admite muchas descripciones válidas y un texto puede parecerse a la referencia sin ser correcto. Las
métricas se agrupan en tres familias: las textuales comparan el texto
generado con uno de referencia, las diagnósticas comprueban si nombra el diagnóstico correcto y las descriptivas
cuentan cuánto contenido clínico menciona. Estas dos últimas no necesitan texto de referencia.

\begin{table}[H]
    \centering
    \caption{Familias de métricas de generación de reportes: qué comparan, qué evalúan y qué no evalúan.}
    \label{tab:marco_familias_texto}
    \small
    \begin{tabularx}{\textwidth}{@{}>{\raggedright\arraybackslash}p{2.3cm}
        >{\raggedright\arraybackslash\hspace{0pt}}X >{\raggedright\arraybackslash\hspace{0pt}}X
        >{\raggedright\arraybackslash\hspace{0pt}}X >{\raggedright\arraybackslash\hspace{0pt}}X@{}}
        \toprule
        \textbf{Familia} & \textbf{Métricas} & \textbf{Compara con} & \textbf{Evalúa} & \textbf{No evalúa} \\
        \midrule
        Textuales & BLEU, ROUGE-L, METEOR, CIDEr, BERTScore & Un texto de referencia
            & Parecido léxico o de significado con esa referencia
            & Si lo descrito está en la imagen \\
        \filasep
        Diagnósticas & Consistencia diagnóstica & La etiqueta de la imagen
            & Si el texto nombra la categoría correcta
            & Si la descripción es fiel a la imagen \\
        \filasep
        Descriptivas & Cobertura de entidades clínicas & Un vocabulario cerrado
            & Cuánto contenido clínico menciona el texto
            & Si lo mencionado es cierto \\
        \bottomrule
    \end{tabularx}
\end{table}

Las métricas textuales, primera familia de la Tabla~\ref{tab:marco_familias_texto}, comparan el texto generado $g$
con un texto de referencia $r$ escrito para la misma imagen. Cuatro de estas métricas cuentan coincidencias de
$n$-gramas, secuencias de $n$ palabras consecutivas, y la quinta compara significados.

\begin{itemize}
    \item \textbf{BLEU.} Calcula, para cada $n$ entre 1 y 4, la precisión $p_n$ de los $n$-gramas de $g$, la
    fracción que también aparece en $r$. Las combina con una media geométrica y multiplica el resultado por una
    penalización de brevedad, porque un texto muy corto puede tener precisión alta sin cubrir la
    referencia~\cite{papineni2002bleu},
    \begin{equation}
        \mathrm{BLEU} = \mathrm{BP}\cdot\exp\!\left(\sum_{n=1}^{4} \tfrac{1}{4}\log p_n\right), \qquad
        \mathrm{BP} =
        \begin{cases}
            1 & \text{si } \ell_g > \ell_r,\\
            e^{\,1 - \ell_r/\ell_g} & \text{si } \ell_g \le \ell_r,
        \end{cases}
        \label{eq:marco_bleu}
    \end{equation}
    donde $\ell_g$ y $\ell_r$ son las longitudes en palabras de los dos textos.

    \item \textbf{ROUGE-L.} Se basa en la subsecuencia común más larga, $\mathrm{LCS}(g, r)$, la secuencia más
    larga de palabras que aparece en ambos textos en el mismo orden. Con ella define una sensibilidad y una
    precisión, y las combina mediante su media armónica~\cite{lin2004rouge},
    \begin{equation}
        R_{\mathrm{lcs}} = \frac{\mathrm{LCS}(g, r)}{\ell_r}, \qquad
        P_{\mathrm{lcs}} = \frac{\mathrm{LCS}(g, r)}{\ell_g}, \qquad
        \text{ROUGE-L} = \frac{2\,R_{\mathrm{lcs}}\,P_{\mathrm{lcs}}}{R_{\mathrm{lcs}} + P_{\mathrm{lcs}}}.
        \label{eq:marco_rouge}
    \end{equation}

    \item \textbf{METEOR.} Alinea las palabras de los dos textos, primero las idénticas, después las que
    comparten la raíz y por último los sinónimos, y sobre las $m$ palabras alineadas calcula la precisión $P$ y
    la sensibilidad $R$, a la que da más peso. Una penalización reduce el puntaje hasta la mitad cuando las
    palabras alineadas forman muchos tramos contiguos, $ch$, en lugar de pocos~\cite{banerjee2005meteor},
    \begin{equation}
        F_{\text{mean}} = \frac{10\,P R}{R + 9P}, \qquad
        \mathrm{Pen} = 0.5\left(\frac{ch}{m}\right)^{3}, \qquad
        \mathrm{METEOR} = F_{\text{mean}}\,(1 - \mathrm{Pen}).
        \label{eq:marco_meteor}
    \end{equation}

    \item \textbf{CIDEr.} Se diseñó para describir imágenes con $J$ referencias por imagen. Representa cada
    texto con un vector $\mathbf{t}^n(\cdot)$ de pesos TF-IDF por $n$-grama, que crecen con la frecuencia del
    $n$-grama en ese texto y decrecen con la cantidad de imágenes en cuyas descripciones aparece, de modo que las
    expresiones comunes pesan poco. Promedia la similitud coseno con las referencias, y después sobre
    $n$~\cite{vedantam2015cider},
    \begin{equation}
        \mathrm{CIDEr}_n = \frac{1}{J}\sum_{j=1}^{J}
        \frac{\mathbf{t}^n(g)^{\top}\mathbf{t}^n(r_j)}{\lVert\mathbf{t}^n(g)\rVert\,\lVert\mathbf{t}^n(r_j)\rVert},
        \qquad
        \mathrm{CIDEr} = \sum_{n=1}^{4} \tfrac{1}{4}\,\mathrm{CIDEr}_n.
        \label{eq:marco_cider}
    \end{equation}

    \item \textbf{BERTScore.} Compara los \emph{embeddings} contextuales de los tokens, obtenidos con un modelo
    de solo codificador como BERT (Sección~\ref{sec:marco_llm_grandes}), en lugar de las palabras, de modo que una
    paráfrasis correcta no se penaliza. Cada token de un texto se empareja con el más parecido del otro por
    similitud coseno~\cite{zhang2020bertscore},
    \begin{equation}
        R_{\mathrm{BERT}} = \frac{1}{|r|}\sum_{r_i \in r}\,\max_{g_j \in g}\, \mathbf{r}_i^{\top}\mathbf{g}_j,
        \qquad
        P_{\mathrm{BERT}} = \frac{1}{|g|}\sum_{g_j \in g}\,\max_{r_i \in r}\, \mathbf{r}_i^{\top}\mathbf{g}_j,
        \label{eq:marco_bertscore}
    \end{equation}
    donde los \emph{embeddings} $\mathbf{r}_i$ y $\mathbf{g}_j$ están normalizados, y el puntaje final
    $F_{\mathrm{BERT}}$ es la media armónica de los dos.
\end{itemize}

Las métricas con referencia suponen que la referencia es la descripción correcta de la imagen, y en dermatología
ese supuesto no está garantizado. Cuando varios dermatólogos marcan por separado los mismos criterios
dermatoscópicos, la concordancia es en su mayoría de pobre a regular, y moderada solo en unos pocos
criterios~\cite{carrera2016validity}. A eso se suma que la referencia de un conjunto puede no haberla escrito un
dermatólogo (Sección~\ref{sec:marco_rep_tarea}). Una descripción generada puede, por lo tanto, diferir de la
referencia sin estar equivocada. Las métricas diagnósticas y descriptivas evitan ese supuesto porque comprueban lo
que el texto afirma.

\begin{itemize}
    \item \textbf{Consistencia diagnóstica.} Mide la fidelidad diagnóstica de la
    Sección~\ref{sec:marco_rep_ejes}: la fracción de los textos que nombran la categoría real de la imagen o un
    sinónimo reconocido,
    \begin{equation}
        \mathrm{FD} = \frac{1}{n}\sum_{i=1}^{n} \mathbf{1}\bigl[\mathrm{dx}(g_i) = y_i\bigr],
        \label{eq:marco_fd}
    \end{equation}
    donde $g_i$ es el texto generado para la imagen $i$, $y_i$ su categoría real y $\mathrm{dx}(\cdot)$ extrae la
    categoría que el texto nombra.

    \item \textbf{Cobertura de entidades clínicas.} Se apoya en una lista cerrada de términos del
    vocabulario estandarizado de la Sección~\ref{sec:marco_rep_tarea}, agrupados en categorías como color, patrón
    o estructura. Es la proporción de esas categorías que el texto menciona,
    \begin{equation}
        \mathrm{cob}(g) = \frac{\bigl|\{k : \exists\, e \in V_k,\ e \text{ aparece en } g\}\bigr|}{K},
        \label{eq:marco_cobertura_ent}
    \end{equation}
    donde $V_k$ es el conjunto de términos de la categoría $k$ y $K$ el número de categorías. Mide cuánto
    contenido clínico describe el texto y no si lo que describe es cierto.
\end{itemize}

Ninguna de estas familias comprueba que lo descrito esté en la imagen, que es lo que define a una alucinación
(Sección~\ref{sec:marco_rep_fallas})~\cite{hartsock2024vision}. Detectarla requiere que un especialista lea el
reporte contra la imagen y lo puntúe según una rúbrica, es decir, según criterios fijados de antemano.

\subsection{Calibración}
\label{sec:marco_eval_calib}

Además de cuánto acierta un modelo, interesa cuánto vale su probabilidad como medida de confianza. La
Sección~\ref{sec:marco_senales_confianza} mostró que un umbral sobre esa probabilidad solo tiene una lectura clara
si el modelo está calibrado, es decir, si sus probabilidades coinciden con sus frecuencias de acierto. El diagrama
de confiabilidad hace visible la calibración: las predicciones se agrupan en $M$ intervalos según su confianza, y
para cada intervalo $B_m$ se compara la exactitud $\mathrm{acc}(B_m)$ con la confianza media
$\mathrm{conf}(B_m)$, que en un modelo calibrado coinciden. El
\textbf{\gls{ece}} promedia la brecha entre las dos pesándola por la cantidad de predicciones de cada intervalo, y
el \textbf{\gls{mce}} toma la peor,
\begin{equation}
    \mathrm{ECE} = \sum_{m=1}^{M} \frac{|B_m|}{n}\,\bigl|\mathrm{acc}(B_m) - \mathrm{conf}(B_m)\bigr|,
    \qquad
    \mathrm{MCE} = \max_{m}\,\bigl|\mathrm{acc}(B_m) - \mathrm{conf}(B_m)\bigr|.
    \label{eq:marco_ece}
\end{equation}
El ECE resume la calibración típica del modelo, y el MCE importa cuando un error de confianza en un intervalo
concreto tiene consecuencias graves~\cite{guo2017calibration}.

% Figura: diagrama de confiabilidad ilustrativo, con datos sintéticos (no son resultados del trabajo).
% Se usa en sec:marco_eval_calib, después de la definición de ECE y MCE.

\begin{figure}[!tbp]
    \centering
    \begin{tikzpicture}[x=6.5mm, y=55mm]

        % Ejes.
        \draw[figtrazo] (0, 0) rectangle (10, 1);
        \foreach \t in {0, 0.2, 0.4, 0.6, 0.8, 1.0} {
            \draw[figtrazo] (\t*10, 0) -- ++(0, -0.02) node[below, font=\footnotesize] {\t};
            \draw[figtrazo] (0, \t) -- ++(-0.12, 0) node[left, font=\footnotesize] {\t};
        }
        \node[font=\small] at (5, -0.14) {confianza de la predicción};
        \node[font=\small, rotate=90] at (-1.35, 0.5) {exactitud};

        % Barras: exactitud por intervalo (sintética, modelo sobreconfiado).
        \foreach \b/\acc in {0/0.04, 1/0.12, 2/0.21, 3/0.28, 4/0.36, 5/0.44, 6/0.53, 7/0.61, 8/0.70, 9/0.82} {
            \pgfmathsetmacro{\medio}{(\b + 0.5)/10}
            \fill[figazul!35] (\b + 0.05, 0) rectangle (\b + 0.95, \acc);
            \draw[figazul!80!black] (\b + 0.05, 0) rectangle (\b + 0.95, \acc);
            % Brecha hasta la diagonal.
            \fill[figbermellon!30, opacity=0.8] (\b + 0.05, \acc) rectangle (\b + 0.95, \medio);
            \draw[figbermellon!80!black, densely dashed] (\b + 0.05, \acc) rectangle (\b + 0.95, \medio);
        }

        % Diagonal de calibración perfecta.
        \draw[figtrazo, thick, dashed] (0, 0) -- (10, 1);

        % Leyenda.
        \fill[figazul!35, draw=figazul!80!black] (0.5, 0.90) rectangle (1.1, 0.95);
        \node[anchor=west, font=\footnotesize] at (1.2, 0.925) {exactitud del intervalo};
        \fill[figbermellon!30, draw=figbermellon!80!black, densely dashed] (0.5, 0.81) rectangle (1.1, 0.86);
        \node[anchor=west, font=\footnotesize] at (1.2, 0.835) {brecha con la confianza};
        \draw[figtrazo, thick, dashed] (0.5, 0.745) -- (1.1, 0.745);
        \node[anchor=west, font=\footnotesize] at (1.2, 0.745) {calibración perfecta};

    \end{tikzpicture}
    \caption{Diagrama de confiabilidad de un modelo sobreconfiado, con datos sintéticos ilustrativos.}
    \label{fig:marco_confiabilidad}
\end{figure}


En el diagrama de la Figura~\ref{fig:marco_confiabilidad}, que corresponde a un modelo sobreconfiado, cada barra es
la exactitud de un intervalo. Las barras quedan por debajo de la diagonal porque el modelo acierta menos de lo que
su confianza indica, y la franja que las separa de ella es la brecha que promedia el ECE.

El \textbf{puntaje de Brier} mide la calibración sin agrupar en intervalos. Es el error cuadrático medio entre
el vector de probabilidades y la codificación de la clase real,
\begin{equation}
    \mathrm{BS} = \frac{1}{n}\sum_{i=1}^{n}\sum_{c=1}^{C} \left(p_{ic} - y_{ic}\right)^2,
    \label{eq:marco_brier}
\end{equation}
donde $p_{ic}$ es la probabilidad que el modelo asigna a la clase $c$ en la imagen $i$, e $y_{ic}$ vale 1 si
$c$ es la clase real de esa imagen y 0 si no~\cite{brier1950verification}. La entropía cruzada de la
Ecuación~\ref{eq:marco_entropia}, promediada sobre el test, se usa con el mismo fin y se llama
\textbf{log-pérdida} (\emph{log loss}). A diferencia del Brier, castiga sin límite asignar una probabilidad cercana a cero a la clase
real. Ninguna de las dos mide solo la calibración, porque bajan tanto cuando las probabilidades se acercan a las
frecuencias reales como cuando el modelo concentra más probabilidad en la clase correcta.

La calibración se puede corregir después del entrenamiento con un conjunto de validación. El
\textbf{\emph{temperature scaling}} divide todos los logits por un único escalar $T>0$ antes de la softmax, como en la
Ecuación~\ref{eq:marco_temperatura}, y ajusta $T$ para minimizar la entropía cruzada en validación. Como no
cambia el orden de los logits, la clase predicha y la exactitud quedan iguales. El \textbf{\emph{vector scaling}}
aprende en cambio un factor y un sesgo para cada clase; corrige más, pero puede cambiar la clase predicha y
con ello la exactitud~\cite{guo2017calibration}.

\subsection{Inferencia estadística}
\label{sec:marco_eval_estadistica}

Todas las métricas anteriores se calculan sobre el test, que es una muestra de las imágenes posibles. Cuando dos
modelos difieren en una o dos centésimas de F1 macro, esa diferencia puede deberse a la composición de imágenes en
el conjunto de test y no a los modelos. Para afirmar que un modelo es mejor que otro hace falta estimar cuánto
variaría la diferencia de una muestra a otra.

El \textbf{\emph{bootstrap}} estima esa variación con los datos disponibles. Se construyen muchas muestras del
mismo tamaño que el test tomando imágenes al azar con reposición, se calcula la métrica en cada una, y la
dispersión de los valores obtenidos aproxima la que tendría la métrica entre muestras
independientes~\cite{efron1979bootstrap}. En la versión pareada, cada remuestreo toma las mismas imágenes para
los dos modelos y registra la diferencia entre ellos. Si el intervalo de percentiles de esas diferencias no
contiene el cero, la diferencia no se explica por la composición del test. Sirve para cualquier métrica
agregada, como el F1 macro, el AUC o la consistencia diagnóstica.

La \textbf{prueba de McNemar} compara los aciertos de dos clasificadores sobre las mismas
imágenes~\cite{mcnemar1947note}. Solo usa los casos discordantes: $b$, los que el primero acierta y el segundo
falla, y $c$, los del caso inverso. Si los dos modelos fueran equivalentes, cada caso discordante tendría la misma
probabilidad de caer de un lado o del otro, y $b$ seguiría una distribución binomial con $b+c$ ensayos y
probabilidad $1/2$. En su versión exacta, el valor $p$ es la probabilidad, bajo esa distribución, de que $b$ y $c$
se separen al menos tanto como en lo observado,
\begin{equation}
    p = \min\!\left(1,\; 2\sum_{i=0}^{\min(b,c)} \binom{b+c}{i} \left(\frac{1}{2}\right)^{b+c}\right).
    \label{eq:marco_mcnemar}
\end{equation}

Cada prueba a nivel $\alpha$ tiene una probabilidad $\alpha$ de declarar una diferencia que no existe. Al comparar
muchos pares de modelos, la probabilidad de que al menos una de las pruebas lo haga crece con el número de pruebas
$m$. La \textbf{corrección de Holm} controla esa probabilidad para el conjunto de las $m$
pruebas~\cite{holm1979simple}. Ordena los valores $p$ de menor a mayor, $p_{(1)} \le \dots \le p_{(m)}$, y compara
cada uno con un umbral que se relaja paso a paso,
\begin{equation}
    p_{(k)} \le \frac{\alpha}{m - k + 1}, \qquad k = 1, \dots, m.
    \label{eq:marco_holm}
\end{equation}
Las hipótesis se rechazan en ese orden mientras se cumpla la desigualdad. En la primera que no la cumple el
procedimiento se detiene, y esa y todas las siguientes se conservan. El primer umbral, $\alpha/m$, es el de la
corrección de Bonferroni, pero los siguientes son menos exigentes, así que Holm rechaza al menos las mismas
hipótesis que Bonferroni con el mismo control.

\subsection{Robustez}
\label{sec:marco_eval_robustez}

La robustez es cuánto cambia una métrica cuando la entrada se perturba de una manera que no altera su contenido
clínico, como la compresión, el ruido, el desenfoque o la reducción de resolución. Permite estimar el efecto de la
variación de adquisición de la Sección~\ref{sec:marco_generalizacion} sin otro conjunto de datos, degradando las
imágenes del test de maneras controladas y con severidad creciente~\cite{hendrycks2019benchmarking}. Sea $v$ el
valor de una métrica sobre el test original y $\tilde{v}$ su valor sobre el test perturbado con una severidad
dada. La caída relativa
\begin{equation}
    \Delta^{\text{rel}} = \frac{v - \tilde{v}}{v}
    \label{eq:marco_robustez}
\end{equation}
mide la sensibilidad del modelo a esa perturbación, y el modelo es más robusto cuanto menor es. Las perturbaciones
sintéticas reproducen solo una parte de la variación real entre equipos y centros, de modo que una caída pequeña
no garantiza robustez ante un cambio de dominio ni reemplaza la validación externa.

\subsection{Interpretabilidad visual}
\label{sec:marco_eval_gradcam}

Las métricas anteriores dicen cuánto acierta un clasificador pero no en qué se basa. \textbf{Grad-CAM}
(\emph{gradient-weighted class activation mapping}) produce un mapa de calor de las regiones de la imagen que
más influyeron en la predicción de una clase. Parte de los mapas de características $A^k$ de la última capa
convolucional, de $i \times j$ posiciones cada uno, y del puntaje $y^c$ de la clase $c$ antes de la softmax.
Mide cuánto importa cada mapa promediando el gradiente del puntaje respecto del mapa, suma los mapas pesados
por esa importancia y aplica la \textbf{ReLU} (\emph{rectified linear unit}), una función de activación que
conserva solo los valores positivos,
\begin{equation}
    \alpha_k^c = \frac{1}{Z}\sum_{i}\sum_{j} \frac{\partial y^c}{\partial A^k_{ij}}, \qquad
    L^c_{\text{Grad-CAM}} = \mathrm{ReLU}\!\left(\sum_k \alpha_k^c A^k\right),
    \label{eq:marco_gradcam}
\end{equation}
donde $Z$ es la cantidad de posiciones del mapa. La ReLU descarta las regiones que bajan el puntaje de la
clase, que suelen corresponder a otras clases~\cite{selvaraju2017gradcam}. El mapa se reescala al tamaño de
la imagen y se superpone a ella. En el Swin, cuya última etapa también devuelve un mapa con disposición
espacial (Sección~\ref{sec:marco_swin}), el cálculo es igual.

El mapa tiene la resolución de la última capa, mucho menor que la de la imagen, y por eso marca regiones
gruesas~\cite{selvaraju2017gradcam}. Que se vea razonable tampoco basta: algunos métodos de saliencia
producen mapas casi iguales con los pesos del modelo aleatorizados, y la inspección visual por sí sola no
distingue un método fiel de uno que no lo es~\cite{adebayo2018sanity}. El mapa muestra dónde se apoyó el
modelo y no si esa región justifica el diagnóstico, pero sirve para detectar un modelo que se apoya en algo
ajeno a la lesión, como los rasgos espurios de la Sección~\ref{sec:marco_generalizacion}.

Con qué datos y con qué protocolo se evaluó un trabajo decide, entonces, si su resultado se sostiene fuera de su
conjunto de origen: si se midió en otra fuente, si se desagregó por grupos y si sus diferencias con otros modelos
se probaron estadísticamente. Ese es el criterio con el que la sección siguiente revisa la literatura.
